\section{\sys Design}
\label{sec:design}
\label{sec:fluxtune}

\subsection{Overview}
\label{sec:ft:overview}
\dhruv{\sout{By now the audience will know what the core challenges are
in solving the problem of fine-tuning LLMs on resource constrained
devices. So have a "preamble" paragraph highlighting the 3 key solution
components and why they interact together. Refer to the overall
diagram. Then end the paragraph with pointing the reader to
sub-sections.} Resolved: the paragraph below names the parts and points
to each subsection. The diagram is still a placeholder.}

\sys keeps the forward-gradient client of \secref{sec:bg:noise}
unchanged and changes the updates and the server side system. (\figref{fig:ft_system}). It
answers each challenge of \secref{sec:background} with one choice. It
updates on a buffer of fixed size $N$, so the signal in each update is
known from count of updates. It steps a fixed fraction $\rho$ of the weight norm,
with $\rho$ as a function of size N (\secref{sec:ft:step}). Because the
steps are weight-relative, the growth they cause is known in closed
form. \dhruv{``known from count of updates'' above should be ``known
from the number of client gradients''. Also: \tabref{tab:ft:notation}
lists $N=50$, $\rho=0.050$, while \secref{sec:ft:step} gives $N=90$
and $\rho_0=0.067$ for DistilBERT; $K$ is used there but never
defined; and \secref{sec:ft:agg} uses ``probes'', ``uploads'',
``aggregation'' and ``pool'', where the terminology says perturbation,
client gradient, server update and buffer.} One rule ties the
choices together. Every constant the server compares against is a
ratio of like quantities: weight norm over weight norm, or count over
count. Such constants do not move with the loss scale, the model or the
stage of training. Algorithm~\ref{alg:fluxtune} states the full loop.

Let $\ttr\in\mathbb{R}^p$ be the trainable weights (the adapters added
to the frozen model), $L$ the loss, and $\norm{\cdot}$ the $\ell_2$
norm. \tabref{tab:ft:notation} lists the notation and the values used.

\begin{table}[t]
\centering
\caption{Notation and values used (AG News, DistilBERT with adapters).}
\label{tab:ft:notation}
\small
\begin{tabular}{@{}lll@{}}
\toprule
Symbol & Meaning & Value \\
\midrule
$p$ & number of trainable weights & 450{,}340 \\
$P$ & perturbations per client gradient & 10 \\
$h$ & finite-difference spacing & 0.01 \\
$u_k$ & client gradient from client $k$ & --- \\
$N$ & buffer size: client gradients per server update & 50 \\
$G$ & buffer sum $\sum_k u_k$ (the step's direction) & --- \\
$C$ & clients training at once & 30 \\
$s$ & scale in $\rho=s\sqrt{PN/p}$ & 1.5 \\
$\rho$ & relative step $\norm{\Delta\ttr}/\norm{\ttr}$ & 0.050 \\
$B$, $\Phi$ & growth budget spent; weight-norm growth $e^B$ & --- \\
$\Phi_{\mathrm{cap}}$ & growth cap (hard stop on $\Phi$) & 3 \\
$a_t$, $m_t$, $g_t$ & accuracy, its 11-eval mean, gain over 150 updates & --- \\
$T_{\mathrm{bins}}$ & shared minibatches, used in turn & 150 \\
\bottomrule
\end{tabular}
\end{table}

\figslot{ft_system}{\sys. Clients send forward-gradient estimates. The
server buffers $N$ of them, steps by a fixed fraction $\rho$ of the
weight norm, tracks the weight-norm growth $\Phi$, and stops on a
plateau or at $\Phi\ge\Phi_{\mathrm{cap}}$.}

% -----------------------------------------------------------------------
\subsection{Signal-calibrated steps}
\label{sec:ft:step}
\label{sec:ft:pool}

The question is how far the server should step on an update that is
mostly noise. \aishwwarya{explain why it is all noise, or refer to the section that explains it is all noise} The step can be sized in gradient units (SGD, with or without a
variance threshold), in per-coordinate weight units (Adam), or relative to the
whole weight vector. \secref{sec:bg:steps} shows that the first two drift as
the weights grow. \sys takes the third.

\myparagraph{A buffer of chosen size}
The server updates once its buffer holds $N$ client gradients, and steps along
their sum $G=\sum_{k=1}^{N}u_k$. By Eq.~\eqref{eq:bg:law}, $N$ fixes how much
of $G$ is signal. A variance threshold controls the same quantity indirectly,
since variance scales as $1/N$ and the threshold's best value drifts with the
model and the weight norm. \sys states it directly, as a count that is fixed
within a run and set once per model and dataset.

\myparagraph{Setting $N$}
The work needed to reach a given accuracy is a fixed number of client
gradients, so $N$ only decides how that work is split into updates: fewer,
larger ones, or more, smaller ones. It is therefore set by system limits, not
tuned for accuracy. A larger $N$ gives a better-aimed step but fewer updates
per hour and more stale arrivals; the step ceiling $\rho\le0.07$ and the
population $n$ bound it from above, and $N\ge\max(C,2K)$ bounds it from below
(staleness $\le1$). We take the largest $N\le\min\!\big(p(0.07/s)^2/P,\;n\big)$
that is a multiple of $K$, which gives $N=90$ for DistilBERT, $100$ for
RoBERTa-large and $70$ for SmolLM2-1.7B. The step in the next paragraph then
follows from $N$.

\myparagraph{A step relative to the weights}
The server discards the length of update $G$ and moves the trainable weights by a
fraction $\rho_t$ of their own norm:
\begin{equation}
\Delta\ttr=-\rho_t\,\norm{\ttr}\frac{G}{\norm{G}},\qquad
\rho_0=s\sqrt{\frac{PN}{p}} .
\label{eq:ft:step}
\end{equation}
A drifting gradient norm leaves the step unchanged (in one run $\norm{u_k}$
rose from $685$ to $1{,}168$ with no effect), and any factor common to all
client gradients cancels, so every client gradient gets weight one. The
trust-ratio form comes from LARS and LAMB~\cite{you2017lars,you2020lamb}.
What is new is how the step is set. The initial step $\rho_0$ scales with the
cosine the buffer can deliver, so a step never outruns its signal by more than
a factor $s$ ($\rho_0=0.067$ for $N=90$). We claim this scaling, not the
constant: $s$ is calibrated once per model (\secref{sec:ft:discussion}).
Because the absolute step still grows with $\norm{\ttr}$ and late updates add
mostly noise, $\rho_t$ decays exponentially from $\rho_0$, which bounds the
total growth of the weight norm.
\dhruv{Consistency check from the intro rewrite: this says $\rho_t$
decays exponentially, but Algorithm~\ref{alg:fluxtune} keeps $\rho$
constant, \secref{sec:ft:budget} says ``at constant $\rho$'', and the
context file's shipped config is $\rho=0.050$ constant. The intro's
C2 claims the absolute step grows with the weights, which holds only
if $\rho_t$ decays slower than $\norm{\ttr}$ grows. Which one ships?}

\myparagraph{Why scale $\rho_0$ with the cosine}
Each update is nearly perpendicular to the weights
($|\cos(\ttr,\Delta\ttr)|<0.0006$), so it grows $\ln\norm{\ttr}$ by about
$\rho_t^2/2$, while its progress along the true gradient is proportional to
$\rho_t\cos(G,\nabla L)$. Growth per unit of progress thus scales as
$\rho_t/\cos$. Choosing $\rho_0\approx s\cos$ holds that ratio at a constant
\weak{derived here, not yet measured}.

\aishwwarya{TODO- add the weight - radian diagram for better explanation}

% -----------------------------------------------------------------------
\subsection{Aggregating averaged client gradients}
\label{sec:ft:agg}

\emph{The question} is what a client should upload and how the server should
combine uploads. A client holds $P$ probes, each a random direction with a
measured slope, and only the server sees several clients. The design space has
two axes: how a client reduces its $P$ probes to one upload (keep the single
best-aligned probe, or combine all of them), and how the server weights the
uploads it sums. Both choices fix how much of $G$ is signal
(Eq.~\eqref{eq:bg:law}); \sys makes the simplest choice on each axis.

\myparagraph{A client uploads the mean of all its probes}
A client draws $P$ Gaussian directions $v_i$ over the trainable slice and
measures the slope of the loss along each with a central difference,
\begin{equation}
d_i=\frac{L(\theta+hv_i)-L(\theta-hv_i)}{2h}\approx\nabla L\!\cdot\! v_i,\qquad
u_k=\frac1P\sum_{i=1}^{P}d_iv_i,\qquad\E[u_k]=\nabla L .
\label{eq:ft:fg}
\end{equation}
This costs $2P$ forward passes and no backward pass. Memory is that of
inference and does not grow with $P$, and the upload is one $p$-vector
whatever $P$ is. Keeping only the probe with the largest $|d_i|$ discards
$P-1$ of the $P$ measurements the client paid for. The mean uses all of them:
its noise power is $\tfrac{p}{P}\norm{\nabla L}^2$, so every added probe
lowers it, and the estimate stays unbiased. In our runs the mean gave
$3.35\times$ more alignment per unit of compute than selection. The argument
applies to any forward-only method that forms its estimate from several
probes per client.

\myparagraph{The server sums uploads with equal weight}
The server sums the $N$ uploads in its buffer, $G=\sum_{k=1}^{N}u_k$, so $G$
is the mean of $NP$ independent probes and
\begin{equation}
\cos(G,\nabla L)\approx\sqrt{\frac{PN}{p}} ,
\label{eq:ft:aim}
\end{equation}
which grows as $\sqrt N$ because signal adds coherently and noise adds as a
random walk. Where the $P$ probes sit does not matter to the server: $P$ probes
in each of $N$ clients and one probe in each of $NP$ clients give the same
$G$. This is what lets \sys trade per-client compute for the number of
clients without changing the algorithm. The argument needs the uploads to be
independent, and they are: the number of independent uploads the pool is
worth, $n_{\mathrm{eff}}$, stayed within $[0.98,1.03]\,N$ across runs. Measured
alignment was about $12\times$ below Eq.~\eqref{eq:ft:aim} (backpropagation
audits on held-out data gave $\cos\approx0.003$), so we use the formula for
its scaling in $(P,N,p)$ and treat the constant as empirical.

Weights on uploads would only matter if some uploads were worth more than
others. We use $\omega_k=1$: because the step discards the length of $G$
(Eq.~\eqref{eq:ft:step}), any factor common to all uploads cancels, and the
remaining freshness- and alignment-based weights changed the aggregate by at
most $0.6\%$ in a measurement \weak{pending the grad-aware comparison}.

\myparagraph{All uploads in a buffer see the same data}
The training data is cut into $T_{\mathrm{bins}}$ fixed local batches, indexed
by \texttt{data\_id}. Aggregation $t$ uses bin $(t-1)\bmod T_{\mathrm{bins}}$
for every upload in its buffer, so the probes being averaged all estimate the
gradient of the same loss. Successive aggregations cycle through the bins.

\myparagraph{Uploads are accepted asynchronously, without staleness correction}
The server keeps $C$ clients running and gives each freed slot the current
weights, so no round waits for the slowest device. An upload computed on
weights $\tau$ aggregations old is added to the buffer unchanged. Since
$\E[\tau]\approx C/N$, keeping $C\le N$ holds staleness near zero: the measured
mean was $0.58$, with $96\%$ of uploads at most one aggregation stale. At this
level a staleness weight had little to correct, so we do not use one.

% -----------------------------------------------------------------------
\subsection{Growth-budgeted stopping}
\label{sec:ft:budget}
\label{sec:ft:stop}

The question is when a run should stop with no target.
\emph{The design space} includes plateau detection on accuracy, a fixed
number of server updates, and a cap on weight growth. Plateau detection
alone has no backstop when evaluation noise hides the plateau
(\secref{sec:bg:unattended}). \sys uses a plateau rule, backed by a cap
on growth that it can compute exactly.

\myparagraph{The growth budget}
Each update is perpendicular to the weights (\secref{sec:bg:growth}).
An update of relative size $\rho_k$ therefore multiplies the weight
norm by exactly $\sqrt{1+\rho_k^2}$. Summing in log space gives the
\emph{growth budget} spent after $t$ server updates:
\begin{equation}
\begin{gathered}
B_t=\tfrac12\sum_{k\le t}\ln\!\left(1+\rho_k^2\right)\approx\tfrac12\sum_{k\le t}\rho_k^2,\\
\Phi_t=e^{B_t}=\norm{\ttr^{(t)}}\big/\norm{\ttr^{(0)}} .
\end{gathered}
\label{eq:ft:phi}
\end{equation}
$B_t$ depends only on the steps taken. It needs no gradients, no
accuracy and no tuned constants. Its prediction matched the measured
norm ratio to within $0.04\%$ (median).

The budget view is what makes $B$ useful. $\sum\rho^2$ is the sum of
squared step sizes, the quantity classic stochastic-approximation
theory requires to stay finite~\cite{robbins1951stochastic}. Here it
also sets how far the weights have grown. So step size and run length
trade off against one budget: larger steps spend it faster and end the
run sooner. At constant $\rho$ the budget reduces to a step counter,
$\Phi_t=(1+\rho^2)^{t/2}$. At $\rho=0.050$, $\Phi$ reaches $3$ after
about $880$ server updates. Any schedule for $\rho$ spends the same
budget differently \tbd{E12}.

\myparagraph{The stop rule}
Let $a_t$ be held-out accuracy after server update $t$, $m_t$ its
11-evaluation trailing mean, and $g_t=(m_t-m_{t-150})/m_t$ its recent
relative gain. Let $n_t$ count the consecutive evaluations with
$g_t\le0.003$. After a warm-up of $450$ server updates, the run stops
when
\begin{equation}
n_t\ge20\quad\lor\quad\Phi_t\ge\Phi_{\mathrm{cap}} .
\label{eq:ft:stop}
\end{equation}
Smoothing and patience keep evaluation noise from ending the run
early. The cap $\Phi_{\mathrm{cap}}=3$ sits below the cliff of
\secref{sec:bg:growth}, so it bounds the damage if the plateau is
missed.

% -----------------------------------------------------------------------
\subsection{Inherited pieces}
\label{sec:ft:inherited}
\label{sec:ft:fg}
\label{sec:ft:async}

Clients run the forward-gradient primitive of \secref{sec:bg:noise},
averaging all $P$ perturbations (Eq.~\eqref{eq:bg:fg}). Server update
$t$ uses shared minibatch $t\bmod T_{\mathrm{bins}}$. All client
gradients in a buffer therefore see the same data, and successive
updates cycle through it. Aggregation is asynchronous, as in
FedBuff~\cite{nguyen2022fedbuff}. The server keeps $C$ clients running
and hands each freed slot the current weights. Client gradients join
the buffer as they arrive, unweighted, whatever their staleness.
Keeping $C\le N$ holds staleness low (\secref{sec:bg:unattended}).

\begin{algorithm}[t]
\caption{\sys}
\label{alg:fluxtune}
\begin{algorithmic}[1]
\Require weights $\theta$; $P,h,N,C,s,\Phi_{\mathrm{cap}},T_{\mathrm{bins}}$
\State $\rho\gets s\sqrt{PN/p}$;\; $B\gets 0$;\; $t\gets 0$;\; buffer $\gets\emptyset$
\While{not stopped}
  \State $b\gets t\bmod T_{\mathrm{bins}}$ \Comment{shared minibatch}
  \State keep $C$ clients busy, each with the current $(\theta,b)$
  \ForAll{client \textbf{in parallel}}
    \State draw $v_1,\dots,v_P\sim\mathcal{N}(0,I_p)$
    \State $d_i\gets[L(\theta{+}hv_i)-L(\theta{-}hv_i)]/2h$ \Comment{on $b$}
    \State send $u\gets\frac1P\sum_i d_iv_i$
  \EndFor
  \State on each arrival, add $u$ to the buffer \Comment{any staleness}
  \If{$|\text{buffer}|\ge N$}
    \State $G\gets\sum_{u\in\text{buffer}}u$
    \If{$\norm{G}>0$}
      \State $\ttr\gets\ttr-\rho\norm{\ttr}\,G/\norm{G}$ \Comment{Eq.~\eqref{eq:ft:step}}
    \EndIf
    \State $B\gets B+\tfrac12\ln(1+\rho^2)$;\; $\Phi\gets e^{B}$ \Comment{Eq.~\eqref{eq:ft:phi}}
    \State $t\gets t+1$;\; buffer $\gets\emptyset$
    \State every two updates, evaluate $a_t$ and update $n_t$
    \If{$n_t\ge20$ \textbf{or} $\Phi\ge\Phi_{\mathrm{cap}}$} \Comment{Eq.~\eqref{eq:ft:stop}}
      \State \textbf{stop}
    \EndIf
  \EndIf
\EndWhile
\end{algorithmic}
\end{algorithm}

% -----------------------------------------------------------------------
\subsection{Limits}
\label{sec:ft:discussion}
All evidence so far is one model (DistilBERT with adapters) on one
dataset (AG News), with one run per arm and evaluation noise of
$\pm0.005$. The constant $s$ is empirical. The measured $\rho/\cos$ was
$14$--$20$, not $1.5$. The cap $\Phi_{\mathrm{cap}}=3$ overshot here.
Accuracy peaked near $\Phi\approx2.6$ and then slid
(\secref{sec:bg:unattended}), and the cap does not transfer between
models. The finite difference is precision-sensitive. With fixed $h$,
the relative perturbation size falls from $0.50$ to about $0.17$ as
$\norm{\ttr}$ triples, which did not hurt training here.

% =============================================================================
% Review comments resolved by the 2026-10-16 rewrite. Their sections moved
% to 02-background.tex or were cut; kept struck through per convention.
% =============================================================================
\dhruv{\sout{if these are papers, they need to be cited} Resolved: LARS
and LAMB are cited.}
\dhruv{\sout{I think you are trying to explain the steps/ intuition here
but its midway between defining the variables and giving the
intuition.} Resolved: the overview gives intuition only; symbols live
in the table.}
\dhruv{\sout{this or theta?} Resolved: $p$ is the count, $\ttr$ the
weights.}
\dhruv{\sout{Sure fwdllm is a baseline but we dont center our solution
around it.} Resolved: ``What FwdLLM does not solve'' is now general
challenges in \secref{sec:background}.}
\dhruv{\sout{Too much detail here and we shouldnt keep it in the main
text as some of it might also be model specific} Resolved: claim-level
only, in \secref{sec:bg:growth}.}
\dhruv{\sout{Its going to be hard for systems people to follow this.
Break it down into two parts.} Resolved: intuition first, then a
skippable ``For the ML reader'' paragraph.}
\dhruv{\sout{It needs to flow a bit better like an explanation/ teaching
of new concepts against previously known concepts.} Resolved: rewritten
as \secref{sec:bg:steps}.}
\aishwwarya{\sout{explain this better - This threshold has to be
fine-tuned for each model, dataset and even varies when we choose log 2
vs log e.} Resolved: \secref{sec:bg:steps}.}
\dhruv{\sout{Bulleted list isn't needed.} Resolved: prose in
\secref{sec:bg:steps}.}
\dhruv{\sout{This section is about our solution, the readers don't care
about fwdllm.} Resolved: design no longer names FwdLLM.}
\aishwwarya{\sout{how do we explain alignment, maybe give context wrt
comparing against backward prop gradients} Resolved: ``cosine with the
true gradient'', contrasted with backprop in \secref{sec:bg:noise}.}
\aishwwarya{\sout{have to explain this better} Resolved: the buffer
paragraph in \secref{sec:ft:step}.}
\aishwwarya{\sout{this is not from the simplified algorithm} Resolved:
the sawtooth schedule is not shipped and was removed.}
\dhruv{\sout{What is this?} Resolved: the ``full controller'' is not
shipped, so its comparison was cut.}
